% Generated draft: requires mathematical and source review.
\documentclass[11pt]{article}
\usepackage[margin=25mm]{geometry}
\usepackage{fontspec}
\setmainfont{Latin Modern Roman}
\usepackage{amsmath,amssymb,mathtools,unicode-math}
\setmathfont{Latin Modern Math}
\usepackage{enumitem,needspace,xurl,hyperref}
\setlength{\parindent}{0pt}
\setlength{\parskip}{5pt}
\setlength{\emergencystretch}{3em}
\newcommand{\N}{\mathbb{N}}
\newcommand{\Z}{\mathbb{Z}}
\newcommand{\Q}{\mathbb{Q}}
\newcommand{\R}{\mathbb{R}}
\newcommand{\C}{\mathbb{C}}
\newcommand{\F}{\mathbb{F}}
\newcommand{\A}{\mathbb{A}}
\newcommand{\PP}{\mathbb{P}}
\newcommand{\E}{\mathbb{E}}
\newcommand{\eps}{\varepsilon}
\newcommand{\dd}{\,\mathrm{d}}
\begin{document}
\section*{UnsolvedMath ID: 30003959}
\textbf{Source problem code:} catalogue entry
\subsection{Definitions and mathematical statement}
Let \(X^{r,\rho}(t)\) denote the random Lorentz-gas particle trajectory named in the supplied record, and let \(\mathcal N\) denote the event described there as ``the particle is not trapped.'' The record does not define the environment distribution, the dynamics, the spatial dimension, the parameters \(r,\rho\), the initial conditions, or the event \(\mathcal N\). These are unresolved parts of the formulation.

For \(T>0\), define the rescaled trajectory by
\[
Y_T(t)=\frac{X^{r,\rho}(Tt)}{\sqrt T},\qquad 0\leq t\leq 1.
\]
The supplied question asks whether, after conditioning on \(\mathcal N\), the laws of \(Y_T\) converge weakly as \(T\to\infty\) to the law of a Wiener process with nondegenerate variance, in the quenched or averaged-quenched setting.

More explicitly, once the missing model data and conditioning laws have been specified, the intended convergence has the schematic form
\[
\mathcal L_{\star}\bigl(Y_T\mid\mathcal N\bigr)
\Longrightarrow \mathcal L(W),
\qquad T\to\infty,
\qquad
\star\in\{\text{quenched},\text{averaged-quenched}\}.
\]
Here \(\mathcal L_{\star}\) denotes the trajectory law in the indicated setting; its precise definition is not supplied. If the spatial state space is \(\R^d\), the usual meaning of a Wiener process with nondegenerate variance is a process \(W\) starting at zero, with continuous paths and independent Gaussian increments satisfying
\[
\mathbb E[W(t)-W(s)]=0,\qquad
\operatorname{Cov}(W(t)-W(s))=(t-s)\Sigma,
\quad 0\leq s\leq t\leq1,
\]
where \(\Sigma\) is a positive definite \(d\times d\) matrix. The dimension and intended covariance convention require confirmation; isotropy is not stated.

This is a statement of the supplied target, subject to unresolved qualifications. The record does not specify a path space or topology for weak convergence, parameter quantifiers, an environment-almost-sure quantifier for the quenched assertion, or the order of averaging and conditioning in the averaged-quenched assertion. The conditioning also requires an appropriate positive-probability nontrapping event, or another explicitly defined conditional law. No complete theorem with these missing hypotheses is asserted.
\subsection{Short English statement}
The supplied problem asks whether a random Lorentz particle, conditioned never to be trapped, has a nondegenerate Brownian limit under the scaling \(X^{r,\rho}(Tt)/\sqrt T\) on \([0,1]\), in a quenched or averaged-quenched formulation. The model, conditioning, parameter ranges, and convergence topology need clarification before this becomes a precise theorem statement. The supplied open status is dated curation, not an independently verified conclusion.
\subsection{Sources}
\begin{itemize}
\item Supplied local UnsolvedMath registry record, permanent ID \(30003959\), title: \emph{Diffusive Limits for the Random Lorentz Gas}; collection: Oberwolfach Reports Open Problems; exact source problem code: \texttt{catalogue entry}. The record gives proposed year 2018 and names no individual proposer.
\item Supplied report citation: \emph{Scaling Limits in Models of Statistical Mechanics} (2019), Oberwolfach Reports, pp.\ 2535--2581, DOI \texttt{10.4171/owr/2018/41}. Source URL and supplied literature-source link labeled ``Oberwolfach Report 41, Lorentz particle'': \url{https://doi.org/10.4171/owr/2018/41}. The supplied background describes Oberwolfach Reports as documenting workshops and their open-problem discussions.
\item Supplied literature-source link labeled ``Lutsko--Toth, Lorentz gas invariance principle'': \url{https://doi.org/10.1007/s00220-020-03852-8}. No further bibliographic details are supplied.
\item Literature assessment dated exactly \texttt{2026-08-21}: the supplied curation labels the problem open and describes the conditioned nontrapped Lorentz-particle invariance principle as a major open problem. It also says that later Lorentz-gas work records a corresponding periodic higher-dimensional invariance problem as open. This is supplied dated curation, not independently verified; its time-sensitive status, cited claims, and applicability to this formulation require independent checking.
\end{itemize}
\end{document}
